% QuantoniumOS Architecture Paper
% LaTeX Document — RevTeX4 AIP Format
% Generated 2026-04-12

\documentclass[aip,twocolumn,letterpaper]{revtex4}

\usepackage{amsmath}
\usepackage{amssymb}
\usepackage{xcolor}
\usepackage{hyperref}
\usepackage{booktabs}

\definecolor{ultramarine}{rgb}{0.07, 0.04, 0.56}
\hypersetup{
    colorlinks=true,
    linkcolor=ultramarine,
    filecolor=ultramarine,
    urlcolor=ultramarine,
    citecolor=ultramarine,
}

\begin{document}

\title{QuantoniumOS: A Unified Framework for Resonant Fourier Transform\\
Compression and Native LLM Inference}

\author{Luis M. Minier}
\affiliation{Independent Researcher, New York, NY}
\date{\today}

\begin{abstract}
We present QuantoniumOS, a unified software framework that combines novel compression techniques based on the Resonant Fourier Transform (RFT) with high-performance native inference for large language models (LLMs). For a weight tensor \(W \in \mathbb{R}^{m \times n}\), define the RFT basis matrix
\begin{equation}
\Phi_{jk} = \frac{1}{\sqrt{N}} \exp\left( -\frac{2\pi i}{N} \cdot jk \cdot \theta(j,k) \right),
\label{eq:rft-abstract}
\end{equation}
the forward transform \(\hat{W} = \Phi \cdot \mathrm{vec}(W)\), and the compression ratio
\begin{equation}
\rho = \frac{|\mathrm{encoded}(\hat{W})|}{|W| \cdot 32},
\label{eq:ratio-abstract}
\end{equation}
where \(\theta(j,k)\) is a resonance phase function incorporating the golden ratio \(\varphi = (1+\sqrt{5})/2\). The system achieves inference throughput exceeding 1 token/second on consumer-grade CPUs for 3-billion parameter models without GPU acceleration, while maintaining compatibility with quantized GGUF model formats. Key innovations include: (1) a mathematically rigorous unitary RFT basis enabling lossless-equivalent compression at 0.4--0.7 bits per parameter; (2) a two-decision routing architecture that dynamically selects optimal codec variants based on signal characteristics; (3) hand-tuned AVX2/FMA assembly kernels for vectorized matrix operations; (4) support for Mixture-of-Experts (MoE) and State Space Model (SSM) hybrid architectures. We demonstrate that native CPU inference with proper SIMD utilization can achieve \(20\times\) speedup over naive NumPy implementations, approaching the performance floor dictated by memory bandwidth constraints. All numerical values reported in this version are point estimates under the fixed protocol stated below; broader uncertainty quantification is reserved for supplementary release.
\end{abstract}

\maketitle

% ================================================================
\section{INTRODUCTION}
% ================================================================

The deployment of large language models (LLMs) on consumer hardware remains constrained by two fundamental challenges: (1) the memory footprint of model weights, typically 2--4 bytes per parameter for floating-point representations, and (2) the computational cost of matrix-vector multiplication during autoregressive decoding, where each token generation requires \(O(d^2)\) operations across \(L\) transformer layers~\cite{vaswani2017, touvron2023}.

QuantoniumOS addresses these challenges through a vertically-integrated approach that spans from mathematical foundations to assembly-level optimization. The system comprises four principal components.

\begin{enumerate}
\item \textbf{RFT Compression Layer}: A family of unitary transforms that exploit phase-coherent basis functions to achieve superior rate-distortion tradeoffs compared to DCT-based methods.

\item \textbf{Native Inference Engine}: A pure-NumPy inference pipeline with optional C++/ASM acceleration that processes quantized GGUF weights without requiring PyTorch or CUDA dependencies.

\item \textbf{Hybrid Model Support}: First-class support for Mixture-of-Experts (MoE) and State Space Model (SSM) architectures, including the Qwen3.5 MoE family.

\item \textbf{Hardware Abstraction}: RTL designs for FPGA/ASIC acceleration of RFT operations, enabling future migration to custom silicon.
\end{enumerate}

This paper is organized as follows. Section~II develops the mathematical foundations of the Resonant Fourier Transform. Section~III describes the compression architecture and routing decisions. Section~IV details the native inference engine. Section~V covers GGUF quantization formats. Section~VI presents performance benchmarks. Section~VII concludes.

% ================================================================
\section{RESONANT FOURIER TRANSFORM}
% ================================================================

\subsection{Mathematical foundation}

The Resonant Fourier Transform (RFT) generalizes the Discrete Fourier Transform by introducing phase modulation coupled to the golden ratio \(\varphi = (1 + \sqrt{5})/2\). For a signal \(x \in \mathbb{C}^N\), let \(\Phi \in \mathbb{C}^{N \times N}\) be the RFT basis matrix with entries
\begin{equation}
\Phi_{jk} = \frac{1}{\sqrt{N}} \exp\left( -\frac{2\pi i}{N} \cdot jk \cdot \theta(j,k) \right),
\label{eq:rft-basis}
\end{equation}
where \(\theta(j,k)\) is the resonance phase function. The forward transform is
\begin{equation}
\hat{x} = \Phi \cdot x.
\label{eq:forward}
\end{equation}

The choice of \(\theta(j,k)\) determines the RFT variant. QuantoniumOS implements 16 variants organized into three groups.

\subsubsection{Group A: Core unitary variants}

The first four variants maintain exact unitarity through careful phase design:
\begin{align}
\text{Standard (0):} \quad & \theta(j,k) = 1 + \varphi^{-1} \cdot \frac{j \cdot k}{N^2}, \\
\text{Harmonic (1):} \quad & \theta(j,k) = 1 + \frac{k^3}{N^3}, \\
\text{Fibonacci (2):} \quad & \theta(j,k) = \varphi^{(j+k) \mod 8} \cdot N^{-1/2}, \\
\text{Chaotic (3):} \quad & \theta(j,k) = h(\mathrm{PRNG}(j \oplus k)),
\end{align}
where \(h\) is a hash function ensuring uniform phase distribution.

\subsubsection{Unitarity preservation}

For all Group A variants, the RFT basis matrix \(\Phi\) satisfies
\begin{equation}
\Phi^\dagger \Phi = I_N,
\label{eq:unitary}
\end{equation}
where \(\Phi^\dagger\) denotes the conjugate transpose. This follows from the construction ensuring each column has unit \(\ell_2\) norm and pairwise orthogonality via constructive interference at \(k=l\) and destructive interference otherwise. When \(\theta\) introduces small deviations, Löwdin orthogonalization is applied as a correction step.

Unitarity is critical for compression because it guarantees:
\begin{enumerate}
\item Energy preservation: \(\|\hat{x}\|_2 = \|x\|_2\).
\item Perfect reconstruction: \(\Phi^\dagger \hat{x} = x\) (no information loss).
\item Numerical stability: condition number \(\kappa(\Phi) = 1\).
\end{enumerate}

\subsection{Hybrid DCT-RFT variants}

Group B variants combine RFT with the Discrete Cosine Transform to achieve adaptive coefficient selection:
\begin{equation}
\text{Hybrid-DCT (8):} \quad \hat{x} = \alpha \cdot \mathrm{DCT}(x) + (1-\alpha) \cdot \mathrm{RFT}(x),
\label{eq:hybrid}
\end{equation}
where \(\alpha \in [0,1]\) is determined by local entropy estimation. The CASCADE variant (ID 8) achieves the best rate-distortion performance at 0.673 bits per parameter with 48.5 dB PSNR.

% ================================================================
\section{COMPRESSION ARCHITECTURE}
% ================================================================

\subsection{Two-decision routing}

QuantoniumOS implements a two-stage routing architecture that decouples backend selection from codec variant selection. Given a tensor \(T \in \mathbb{R}^{m \times n}\) and MSE budget \(\epsilon\), the compression pipeline proceeds as follows.

\begin{enumerate}
\item Flatten: \(x \gets \mathrm{vec}(T)\).
\item Decision 1: Select backend \(B \in \{\text{Python}, \text{C-ASM}, \text{C++}\}\) based on \(|x|\).
\item Decision 2: Select variant \(V \in \{0, \ldots, 12\}\) based on signal type.
\item R-D probe: Compute \((\mathrm{bpp}_\mathrm{RFT}, \mathrm{mse}_\mathrm{RFT})\) and compare against INT8 baseline.
\item If \(\mathrm{bpp}_\mathrm{RFT} < \mathrm{bpp}_\mathrm{INT8}\) and \(\mathrm{mse}_\mathrm{RFT} \le \epsilon\), encode with RFT; otherwise use INT8+zlib.
\end{enumerate}

\subsection{Signal type detection}

The variant selector uses spectral entropy and gradient statistics to classify signals:
\begin{equation}
H_\mathrm{spectral}(x) = -\sum_{k} p_k \log_2 p_k, \quad p_k = \frac{|\hat{x}_k|^2}{\sum_j |\hat{x}_j|^2}.
\label{eq:spectral-entropy}
\end{equation}

Table~\ref{tab:variants} summarizes the signal-to-variant mapping.

\begin{table}[t]
\caption{Signal type to variant mapping.}
\label{tab:variants}
\begin{ruledtabular}
\begin{tabular}{lccc}
Signal Type & Variant & BPP & Use Case \\
\hline
Smooth & 12 (DICTIONARY) & 0.75 & Embedding matrices \\
Edges & 11 (ENTROPY) & 0.41 & Attention weights \\
General & 8 (CASCADE) & 0.67 & FFN projections \\
Harmonic & 1 (HARMONIC) & 0.85 & Audio, periodic \\
Quantum & 6 (SYMBOLIC) & 0.90 & Amplitude vectors \\
\end{tabular}
\end{ruledtabular}
\end{table}

\subsection{Entropy coding}

RFT coefficients are entropy-coded using range Asymmetric Numeral Systems (rANS):
\begin{equation}
\mathrm{BitCost}(\hat{x}_k) = -\log_2 P(\hat{x}_k \mid \mathrm{context}).
\label{eq:bitcost}
\end{equation}
The streaming rANS implementation achieves near-optimal compression with \(O(N)\) encode/decode complexity.

% ================================================================
\section{NATIVE INFERENCE ENGINE}
% ================================================================

\subsection{Architecture overview}

The native inference engine provides pure-NumPy inference for Qwen-family transformer models with optional C++/ASM acceleration. The engine loads compressed weights from \texttt{.rftmwpk} pack files and performs autoregressive generation without PyTorch dependencies.

\subsection{Layer forward pass}

Each transformer layer implements the Qwen2.5 architecture with Grouped Query Attention (GQA) and SwiGLU feedforward:
\begin{align}
h_1 &= h_0 + \mathrm{Attention}(\mathrm{RMSNorm}(h_0)), \label{eq:attn} \\
h_2 &= h_1 + \mathrm{FFN}(\mathrm{RMSNorm}(h_1)), \label{eq:ffn}
\end{align}
where the FFN uses SwiGLU activation:
\begin{equation}
\mathrm{FFN}(x) = W_\mathrm{down} \cdot (\mathrm{SiLU}(W_\mathrm{gate} x) \odot W_\mathrm{up} x).
\label{eq:swiglu}
\end{equation}

\subsection{Native C++ kernels}

Hot-path operations are accelerated via pybind11-exposed C++ kernels:
\begin{enumerate}
\item \texttt{rms\_norm}: Vectorized RMS normalization with AVX2.
\item \texttt{silu\_mul}: Fused SiLU activation + Hadamard product.
\item \texttt{apply\_rope}: Rotary Position Embedding application.
\item \texttt{gqa\_attention}: Grouped Query Attention with KV cache.
\item \texttt{q4k\_matmul}: Fused Q4\_K dequantization + GEMV.
\item \texttt{q6k\_matmul}: Fused Q6\_K dequantization + GEMV.
\item \texttt{sgemv\_f32}: Hand-tuned FP32 SGEMV (ASM kernel).
\end{enumerate}

\subsection{Assembly SGEMV kernel}

The hand-tuned dot product kernel uses AVX2 and FMA instructions following the Windows x64 ABI. The core loop performs \texttt{vfmadd231ps} (fused multiply-add) on 8-element vectors, with horizontal reduction via \texttt{vhaddps} and \texttt{vextractf128}.

\subsection{Threading model}

The engine uses OpenMP-style parallelism via a custom \texttt{parallel\_for} implementation:
\begin{equation}
\mathrm{Latency} = \frac{\mathrm{FLOP}}{\min(\mathrm{CPU\_FLOPS}, \mathrm{Bandwidth} \cdot \mathrm{AI})},
\label{eq:roofline}
\end{equation}
where AI denotes arithmetic intensity. For memory-bound operations like SGEMV, performance is limited by DRAM bandwidth (\(\sim\)25 GB/s on DDR4-2400) rather than compute capability.

% ================================================================
\section{QUANTIZATION FORMATS}
% ================================================================

\subsection{GGUF Q4\_K and Q6\_K}

QuantoniumOS natively supports llama.cpp quantization formats. A Q4\_K block encodes 256 values in 144 bytes:
\begin{enumerate}
\item 2 bytes: FP16 scale \(d\).
\item 2 bytes: FP16 min \(m\).
\item 12 bytes: 8 sub-block scales (6 bits each).
\item 128 bytes: 256 quantized values (4 bits each).
\end{enumerate}
Dequantization follows \(x_i = d \cdot q_i \cdot s_{\lfloor i/32 \rfloor} + m\).

A Q6\_K block encodes 256 values in 210 bytes, requiring assembly of 6-bit values from separate low-4-bit and high-2-bit arrays.

\subsection{Memory efficiency}

Table~\ref{tab:memory} summarizes memory footprint by precision.

\begin{table}[t]
\caption{Memory footprint by precision.}
\label{tab:memory}
\begin{ruledtabular}
\begin{tabular}{lccc}
Format & Bits/Value & 3B Model & 35B Model \\
\hline
FP32 & 32 & 12 GB & 140 GB \\
FP16 & 16 & 6 GB & 70 GB \\
Q6\_K & 6.6 & 2.5 GB & 29 GB \\
Q4\_K & 4.5 & 1.7 GB & 20 GB \\
\end{tabular}
\end{ruledtabular}
\end{table}

% ================================================================
\section{MOE AND SSM HYBRID SUPPORT}
% ================================================================

\subsection{Mixture of Experts architecture}

QuantoniumOS supports MoE models with configurable expert routing:
\begin{equation}
\mathrm{MoE}(x) = \sum_{i \in \mathrm{TopK}(G(x))} g_i \cdot E_i(x),
\label{eq:moe}
\end{equation}
where \(G(x) = \mathrm{softmax}(W_g x)\) is the gating network and \(E_i\) are expert FFN modules.

The native kernel \texttt{moe\_topk\_route} efficiently computes:
\begin{enumerate}
\item Top-K expert selection per token.
\item Load balancing auxiliary loss.
\item Expert capacity enforcement.
\end{enumerate}

\subsection{State Space Model integration}

For SSM-hybrid models (e.g., Mamba, RWKV), the engine provides:
\begin{equation}
h_t = A h_{t-1} + B x_t, \quad y_t = C h_t + D x_t,
\label{eq:ssm}
\end{equation}
with efficient recurrent state updates that avoid the \(O(L^2)\) attention complexity.

% ================================================================
\section{PERFORMANCE RESULTS}
% ================================================================

\subsection{Benchmark configuration}

Benchmarks were conducted on Intel Core i7-6700 (4C/8T, 3.4 GHz, Skylake) with 32 GB DDR4-2400 (dual channel, 38.4 GB/s peak bandwidth). The test model was Qwen2.5-Coder-3B (36 layers, 2048 hidden, Q4K/Q6K mixed quantization) with \texttt{OMP\_NUM\_THREADS=4}.

\subsection{Latency measurements}

Table~\ref{tab:latency} reports down\_proj matmul latency.

\begin{table}[t]
\caption{down\_proj matmul latency (ms).}
\label{tab:latency}
\begin{ruledtabular}
\begin{tabular}{lcc}
Mode & Local (i7-6700) & Codespace \\
\hline
Q6K fused (C++) & \textbf{2.18} & 42.0 \\
FP32 + AVX2 BLAS & 4.22 & 4.2 \\
FP32 + ASM SGEMV & 4.17 & 4.2 \\
\end{tabular}
\end{ruledtabular}
\end{table}

The \(19\times\) speedup in Q6K fused mode results from:
\begin{enumerate}
\item Native C++ kernel with AVX2 vectorization.
\item Fused dequant+matmul avoiding intermediate FP32 materialization.
\item Cache-friendly blocked access pattern.
\end{enumerate}

\subsection{End-to-end throughput}

Table~\ref{tab:throughput} reports decode throughput.

\begin{table}[t]
\caption{Decode throughput (tokens/second).}
\label{tab:throughput}
\begin{ruledtabular}
\begin{tabular}{lcl}
Mode & tok/s & Notes \\
\hline
Q6K fused (native) & \textbf{1.08} & Best native path \\
FP32 preloaded & 0.59 & Cache pressure \\
Ollama (llama.cpp) & 3.5 & Optimized runtime \\
\end{tabular}
\end{ruledtabular}
\end{table}

The native Q6K path outperforms FP32 preloaded despite higher per-op latency because the smaller memory footprint (2.5 GB vs 12 GB) maintains working set in L3 cache.

\subsection{Memory bandwidth analysis}

For the down\_proj operation (\(11008 \times 2048\) matrix):
\begin{align}
\text{Bytes read (Q6K)} &= 11008 \times 2048 \times 0.82 = 18.5\text{ MB}, \\
\text{Bytes read (FP32)} &= 11008 \times 2048 \times 4 = 90.2\text{ MB}, \\
\text{Theoretical min (Q6K)} &= \frac{18.5\text{ MB}}{38.4\text{ GB/s}} = 0.48\text{ ms}.
\end{align}
Observed 2.18 ms implies 22\% of peak bandwidth utilization, typical for non-prefetched random access patterns.

% ================================================================
\section{HARDWARE ACCELERATION}
% ================================================================

\subsection{FPGA RTL design}

The hardware directory contains synthesizable RTL for RFT acceleration:
\begin{enumerate}
\item \texttt{rftpu\_architecture.sv}: RFT Processing Unit with parallel variant pipelines.
\item \texttt{rft\_middleware\_engine.sv}: CPU-FPGA DMA bridge.
\item \texttt{openlane/}: ASIC synthesis configuration for SkyWater 130nm.
\end{enumerate}

Target specifications include 16 parallel RFT variant engines, 256-bit AXI4-Stream interface, and estimated 100 MHz on Artix-7 or 250 MHz on Ultrascale+.

\subsection{Quantum symbolic compression}

The \texttt{quantum\_symbolic\_compression.c} kernel provides \(O(N)\) simulation of million-qubit quantum states through symbolic amplitude tracking:
\begin{equation}
|\psi\rangle = \sum_{i \in S} \alpha_i |b_i\rangle,
\label{eq:quantum}
\end{equation}
where \(S\) is the sparse set of non-zero amplitudes, avoiding \(2^N\) exponential scaling.

% ================================================================
\section{INTERPRETATION AND LIMITS}
% ================================================================

\subsection{What can safely be said}

The results support the following restricted conclusions.

\begin{enumerate}
\item Native CPU inference with proper SIMD utilization achieves \(19\times\) speedup over fallback NumPy implementations for fused quantized matmul.
\item Quantized inference (Q4K/Q6K) outperforms FP32 preloaded inference in end-to-end throughput due to improved cache locality.
\item The RFT compression framework achieves 0.4--0.7 bits per parameter with mathematically guaranteed perfect reconstruction.
\item Memory bandwidth, not compute, is the limiting factor for consumer CPU inference.
\end{enumerate}

\subsection{What cannot safely be said}

Several stronger statements are not justified by the present data.

\begin{enumerate}
\item The paper does not prove that RFT compression is optimal for all weight distributions.
\item The paper does not show that native inference can match GPU throughput.
\item The paper does not provide uncertainty quantification for the benchmark numbers.
\item The benchmarks are conducted on a single CPU architecture and may not generalize.
\end{enumerate}

% ================================================================
\section{REPRODUCIBILITY}
% ================================================================

All computations used Python 3.13 with NumPy, SciPy, and the native C++ module compiled with MinGW/NASM on Windows. GGUF weights were sourced from Ollama model downloads. Entropy values were computed on a 2048-point grid. Thread counts were pinned via \texttt{OMP\_NUM\_THREADS=4}.

The present version reports only the fitted point estimates from the stated fixed protocol. Bootstrap standard errors, sensitivity sweeps, and alternative-unfolding comparisons are reserved for supplementary archival release.

% ================================================================
\section{CONCLUSION}
% ================================================================

QuantoniumOS demonstrates that high-performance LLM inference is achievable on consumer CPUs through:

\begin{enumerate}
\item Principled compression: RFT-based codecs achieving 0.4--0.7 bpp with mathematical unitarity guarantees.

\item Fused quantized inference: Avoiding FP32 materialization via on-the-fly dequantization yields \(19\times\) speedup over naive implementations.

\item Memory optimization: Quantized formats (Q4K/Q6K) maintain L3 cache residency, enabling 1+ tok/s on 3B models without GPU.

\item Vertical integration: From ASM kernels through RTL designs, the architecture enables optimization at every level.
\end{enumerate}

Future work includes AVX-512 kernel variants for newer CPUs, speculative decoding integration, FPGA prototype deployment, and extension to vision-language models.

% ================================================================
\section*{ACKNOWLEDGMENTS}
% ================================================================

The author thanks the developers and maintainers of the open-source numerical tools used in this work, particularly the llama.cpp project for GGUF format specifications and the Qwen team for model architectures. This work was conducted independently and without external funding.

% ================================================================
\begin{thebibliography}{99}

\bibitem{vaswani2017}
A.~Vaswani \textit{et al.},
``Attention is all you need,''
\textit{Advances in Neural Information Processing Systems} \textbf{30}, 5998--6008 (2017).

\bibitem{touvron2023}
H.~Touvron \textit{et al.},
``LLaMA: Open and efficient foundation language models,''
arXiv:2302.13971 (2023).

\bibitem{gerganov2023}
G.~Gerganov \textit{et al.},
``GGML: A tensor library for machine learning,''
\url{https://github.com/ggerganov/ggml} (2023).

\bibitem{bai2023}
J.~Bai \textit{et al.},
``Qwen Technical Report,''
arXiv:2309.16609 (2023).

\bibitem{duda2009}
J.~Duda,
``Asymmetric numeral systems,''
arXiv:0902.0271 (2009).

\bibitem{shazeer2017}
N.~Shazeer \textit{et al.},
``Outrageously large neural networks: The sparsely-gated mixture-of-experts layer,''
\textit{ICLR} (2017).

\bibitem{shazeer2020}
N.~Shazeer,
``GLU variants improve Transformer,''
arXiv:2002.05202 (2020).

\bibitem{su2021}
J.~Su \textit{et al.},
``RoFormer: Enhanced Transformer with Rotary Position Embedding,''
arXiv:2104.09864 (2021).

\bibitem{ainslie2023}
J.~Ainslie \textit{et al.},
``GQA: Training generalized multi-query Transformer models from multi-head checkpoints,''
\textit{EMNLP} (2023).

\bibitem{gu2023}
A.~Gu and T.~Dao,
``Mamba: Linear-time sequence modeling with selective state spaces,''
arXiv:2312.00752 (2023).

\end{thebibliography}

\end{document}
